\documentclass[10pt,a4paper]{amsart}
\usepackage[margin=19mm]{geometry}
\usepackage{amsmath,amssymb,mathtools}
\usepackage[scr=rsfs,cal=euler]{mathalfa}
\usepackage{xfrac}
\usepackage[unicode,hidelinks,bookmarks=false]{hyperref}
\newcommand{\ie}{i.e.\ }
\newcommand{\cf}{cf.\ }
\newcommand{\resp}{respectively\ }
\newcommand{\Subsection}{\subsection}
\providecommand{\nobreakdash}{\nobreak\hbox{-}}
\setlength{\emergencystretch}{2em}
\multlinegap0pt
\usepackage{amssymb}
\usepackage{mathtools}
\newtheorem{theorem}{Theorem}
\title{\bfseries Period Relations for Theta Products and Elliptic Integral Moments}
\author{Dianbin Bao}
\date{Source: arXiv:2608.26514v2 (CC BY 4.0)}
\begin{document}
\maketitle

\noindent\emph{Source and licence.} \bfseries Period Relations for Theta Products and Elliptic Integral Moments. Version arXiv:2608.26514v2 (CC BY 4.0), distributed by arXiv under the Creative Commons Attribution 4.0 International licence, http://creativecommons.org/licenses/by/4.0/, as recorded in the arXiv licence metadata for this exact version and shown on the abstract page https://arxiv.org/abs/2608.26514v2. Full licence notice: \texttt{licenses/arxiv-2608.26514-CC-BY-4.0.txt}.\par\smallskip

\noindent\textit{Editorial scope.} Counted unit: the article's theorem theorem (source0.tex lines 914--923). The article's complete original proof of it is delivered with it (source0.tex lines 925--1159, one \texttt{proof} environment of 235 lines). No mathematics is written, repaired or re-derived: the statement and the proof are the article's own lines, in the article's own notation, and no retained line is substituted or re-flowed.  No further source-local context is needed: every cross-reference of every retained line resolves inside the two delivered spans.  Nothing was omitted from any retained span, so the package carries no omission record.  No retained line contains a citation, so the wrapper carries no bibliography and no bibliography entry.  No retained line contains a figure environment, an image-inclusion command or drawing code, so no image is delivered and the package declares no asset dependency.  Editorial formatting only, in addition: an AMS \texttt{amsart} preamble that restates the article's own declarations and macros used by the retained lines, namely the environments \texttt{theorem} on separate counters, together with the standard packages those lines need.  The retained environments print in this document as Theorem 1 in order of appearance, whereas the article prints the same items with the numbering of its own full document.  The complete native source of the article as distributed in the arXiv e-print for this exact version is bundled unchanged as \texttt{sources/208639/source0.tex}.
\begin{theorem}
	The period polynomial associated with $G_0(\tau)=g(\tau/2)$ is
	\[
	P_g(X)=\frac{P_g(0)}{2}\left(2+5iX-5X^2-2iX^3\right).
	\]
	Equivalently,
	\[
	r_g(X)=r_g(0)\left(1+5iX-10X^2-8iX^3\right).
	\]
\end{theorem}

\begin{proof}
	We first determine the action of $S$ on $G_0$ directly from the transformation law of the Dedekind eta function.  Since
	\[
	G_0(\tau)=\eta(\tau/2)^4\eta(\tau)^2\eta(2\tau)^4,
	\]
	we have
	\[
	G_0\left(-\frac{1}{\tau}\right)
	=
	\eta\left(-\frac{1}{2\tau}\right)^4
	\eta\left(-\frac{1}{\tau}\right)^2
	\eta\left(-\frac{2}{\tau}\right)^4.
	\]
	Using
	\[
	\eta\left(-\frac{1}{\tau}\right)=\sqrt{-i\tau}\,\eta(\tau),
	\]
	in the three forms
	\[
	\eta\left(-\frac{1}{2\tau}\right)=\sqrt{-2i\tau}\,\eta(2\tau),
	\qquad
	\eta\left(-\frac{1}{\tau}\right)=\sqrt{-i\tau}\,\eta(\tau),
	\]
	and
	\[
	\eta\left(-\frac{2}{\tau}\right)
	=
	\sqrt{-\frac{i\tau}{2}}\,\eta\left(\frac{\tau}{2}\right),
	\]
	we obtain
	\[
	G_0\left(-\frac{1}{\tau}\right)
	=
	(-2i\tau)^2(-i\tau)\left(-\frac{i\tau}{2}\right)^2
	G_0(\tau)
	=
	(-i\tau)^5G_0(\tau).
	\]
	Thus
	\[
	G_0|_5S=-iG_0.
	\]
	
	Substituting $\tau=-1/z$ in the period integral therefore gives
	\[
	P_0(X)=-iX^3P_0\left(-\frac{1}{X}\right).
	\]
	Write
	\[
	P_0(X)=a+bX+cX^2+dX^3.
	\]
	Comparison of coefficients gives
	\[
	c=ib,\qquad d=-ia,
	\]
	and hence
	\[
	P_0(X)=a+bX+ibX^2-iaX^3.
	\]
	
	We next use the Gaussian theta series representation.  Since
	\[
	e^{\pi i(n^2+m^2)}=(-1)^{n^2+m^2}=(-1)^{n+m},
	\]
	translation by $1$ gives
	\[
	G_0(\tau+1)=G_1(\tau).
	\]
	Similarly,
	\[
	G_1(\tau+1)=G_0(\tau).
	\]
	
The Gaussian theta-series representation also yields
\[
G_0(\tau)+G_1(\tau)=-8G_0(2\tau).
\]
Indeed, $1+(-1)^{m+n}$ restricts the Gaussian sum to the index two
sublattice $(1+i)\mathbb Z[i]\subset\mathbb Z[i]$; under the corresponding
change of variables the norm is multiplied by $2$ and the quartic factor by
$(1-i)^4=-4$, giving the stated identity.	
	
	
	
	
%	The same lattice representation also gives an additional level-$2$ relation.  Indeed,
%	\[
%	G_0(\tau)+G_1(\tau)
%	=
%	\frac{1}{4}\sum_{m,n\in\mathbb Z}
%	\left(1+(-1)^{m+n}\right)
%	(n-im)^4e^{\pi i\tau(n^2+m^2)}.
%	\]
%	Only the terms with $m+n$ even remain, so
%	\[
%	G_0(\tau)+G_1(\tau)
%	=
%	\frac{1}{2}
%	\sum_{\substack{m,n\in\mathbb Z\\m+n\ {\rm even}}}
%	(n-im)^4e^{\pi i\tau(n^2+m^2)}.
%	\]
%	Put
%	\[
%	r=\frac{n+m}{2},\qquad s=\frac{n-m}{2}.
%	\]
%	Then
%	\[
%	n=r+s,\qquad m=r-s,
%	\]
%	and
%	\[
%	n^2+m^2=2(r^2+s^2).
%	\]
%	Moreover,
%	\[
%	n-im=(1-i)(r+is),
%	\]
%	and since $(1-i)^4=-4$,
%	\[
%	(n-im)^4=-4(r+is)^4.
%	\]
%	It follows that
%	\[
%	G_0(\tau)+G_1(\tau)
%	=
%	-2\sum_{r,s\in\mathbb Z}
%	(r+is)^4e^{2\pi i\tau(r^2+s^2)}.
%	\]
%	Replacing $s$ by $-s$ in the sum gives
%	\[
%	\sum_{r,s\in\mathbb Z}
%	(r+is)^4e^{2\pi i\tau(r^2+s^2)}
%	=
%	\sum_{r,s\in\mathbb Z}
%	(r-is)^4e^{2\pi i\tau(r^2+s^2)}.
%	\]
%	By the definition of $G_0$,
%	\[
%	G_0(2\tau)
%	=
%	\frac{1}{4}\sum_{r,s\in\mathbb Z}
%	(r-is)^4e^{2\pi i\tau(r^2+s^2)},
%	\]
%	and therefore
%	\[
%	G_0(\tau)+G_1(\tau)=-8G_0(2\tau).
%	\]
%	
Passing to period polynomials therefore gives
\[
P_1(X)=-P_0(X)-\frac12P_0(2X).
\]
	Since
	\[
	P_0(X)=a+bX+ibX^2-iaX^3,
	\]
	this gives
	\[
	P_1(X)
	=
	-\frac{3}{2}a-2bX-3ibX^2+5iaX^3.
	\]
	
	It remains to determine the ratio $b/a$.  We apply the same ideal-triangle argument used above, integrating
	\[
	G_0(\tau)(\tau-X)^3\,d\tau
	\]
	around the positively oriented triangle with vertices $0$, $i\infty$, and $-1$.
	
	The side from $0$ to $i\infty$ contributes
	\[
	P_0(X).
	\]
	On the side from $i\infty$ to $-1$, put $\tau=z-1$.  Since
	\[
	G_0(z-1)=G_1(z),
	\]
	this side contributes
	\[
	-P_1(X+1).
	\]
	For the third side, put
	\[
	\tau=-\frac{1}{z+1}.
	\]
	Using the $S$- and $T$-relations,
	\[
	G_0\left(-\frac{1}{z+1}\right)
	=
	(-i(z+1))^5G_0(z+1)
	=
	-i(z+1)^5G_1(z).
	\]
	Therefore the third side contributes
	\[
	iX^3P_1\left(-1-\frac{1}{X}\right).
	\]
	Cauchy's theorem gives the Manin relation
	\[
	P_0(X)-P_1(X+1)
	+iX^3P_1\left(-1-\frac{1}{X}\right)=0.
	\]
	
Taking the constant term in the Manin relation gives
\[
a+\left(\frac{3}{2}a+2b+3ib-5ia\right)+5a=0,
\]
and hence
\[
b=\frac{5i}{2}a.
\]
Therefore
\[
P_0(X)=a\left(1+\frac{5i}{2}X-\frac{5}{2}X^2-iX^3\right).
\]
Since $a=P_0(0)=P_g(0)$, we obtain
\[
P_g(X)=\frac{P_g(0)}{2}\left(2+5iX-5X^2-2iX^3\right).
\]
	
	Finally,
	\[
	P_g(X)=16r_g\left(\frac{X}{2}\right),
	\]
	so
	\[
	r_g(X)=\frac{1}{16}P_g(2X).
	\]
	Since $P_g(0)=16r_g(0)$, we obtain
	\[
	r_g(X)
	=
	r_g(0)\left(1+5iX-10X^2-8iX^3\right).
	\]
\end{proof}

\end{document}
